\chapter{Presentation, Peer Review, and Final Research Package}

\WeekHeader{10}{Presentation, Peer Review, and Final Research Package}{Communicate the project clearly and leave behind a reusable research package.}{Final report, presentation, clean project folder, and continuation note.}

\section{Why This Matters}
The final week is about making the work useful to others. A strong undergraduate project should be understandable to a future student, the dsDNA Core team, a collaborator, or a reviewer. The final package should preserve what was done, why it was done, and how to continue.

\section{Presentation Structure}
\begin{enumerate}
  \item Research question and system.
  \item Why the chemical evidence matters.
  \item Data source and uafR workflow.
  \item Main figure or result.
  \item Evidence supporting interpretation.
  \item Limitations.
  \item Next experiment or analysis.
\end{enumerate}

\section{Peer Review}
Students review each other's final products using the rubric appendix. Reviewers should focus on:
\begin{itemize}
  \item Clarity of research question.
  \item Traceability of chemical evidence.
  \item Quality of cleaned variables and matrices.
  \item Figure readability.
  \item Strength of interpretation.
  \item Reproducibility.
\end{itemize}

\section{Final Package Checklist}
\begin{checklistbox}{Required Handoff Contents}
\begin{itemize}
  \item \path{README.md} explains the project and how to rerun the analysis.
  \item \path{data_raw/} contains unchanged input files or instructions for secure access.
  \item \path{data_processed/} contains derived tables.
  \item \path{scripts/} contains numbered scripts.
  \item \path{figures/} contains final figures.
  \item \path{results/} contains saved enrichment objects and summary tables.
  \item \path{reports/} contains final report and presentation.
  \item \path{logs/} contains research log and session information.
\end{itemize}
\end{checklistbox}

\section{Independent Extension}
Students prepare a continuation note:
\begin{itemize}
  \item What worked.
  \item What failed.
  \item Which chemicals or traits deserve follow-up.
  \item Which database fields were most useful.
  \item Which results should not be overinterpreted.
  \item What the next student should do first.
\end{itemize}

\section{Deliverables}
\begin{tabularx}{\textwidth}{p{0.25\textwidth}YY}
\DeliverableTableHeader
Final report & Revised after dsDNA Core and peer feedback. & PDF\\
Presentation & 8 to 12 minute talk or poster-style slide deck. & PDF or slides\\
Research package & Clean folder with scripts, data, figures, results, and logs. & Project folder\\
Handoff note & Explains next steps for the project. & Markdown or PDF\\
\end{tabularx}

\section{Quality Checklist}
\begin{checklistbox}{Program Completion Standard}
\begin{itemize}
  \item The dsDNA Core team can rerun or inspect the analysis without guessing file locations.
  \item The final interpretation is evidence-based and appropriately cautious.
  \item The student can answer where each major result came from.
  \item The project creates a foundation for continued research.
\end{itemize}
\end{checklistbox}
